\subsection{从一小柱液体看压强随高度的变化}
压强不是向下的力，而是作用在所选小面上的法向力除以面积。液体内部同一点的压强是标量，受压面的朝向不同，压力方向随之改变。因此不能因为重力向下，就把侧壁受到的液体压力也画成向下。

取截面积为 $S$、底面高度为 $z$、顶面高度为 $z+\dd z$ 的竖直小柱。周围液体在底面给它向上的力 $p(z)S$，在顶面给它向下的力 $p(z+\dd z)S$；侧面压力没有竖直分量。向上为正，平衡式为
\[
 p(z)S-p(z+\dd z)S-\rho gS\dd z=0.
\]
减去相同项、除以 $S\dd z$ 并取极限，才得到 $\dd p/\dd z=-\rho g$。负号表示向上压强降低；若改用向下的深度 $h$，则 $\dd p/\dd h=\rho g$。对密度恒定的液体，从自由面 $p=p_a$ 积分到深度 $h$，有
\[
 p(h)-p_a=\int_0^h\rho g\,\dd h'=\rho gh.
\]
$p$ 是绝对压强，$p-p_a$ 是相对大气的表压。自由面与大气相通并不意味着 $p=0$，而是表压为零。比较受压面两侧的合力时，才可能把相同的大气压贡献抵消。

\subsection{不同液层与受压面的合力}
密度分层时，应在每层分别累加压强增量。在同一连通静止液体中取水平小柱，水平外体积力为零，压强差必须为零；这就是等高处等压的力学依据。若两点属于互不连通的液体，或者有水平加速度，就不能仅凭高度相同认定压强相等。

\begin{origexamplebox}{两层液体中的侧壁压力}\label{ch13:ex:teaching-layered}
一个敞口容器中，上层液体密度为 $\rho_1$、厚度为 $h_1$，下层密度为 $\rho_2>\rho_1$、厚度为 $h_2$。竖直侧壁宽度为 $b$，外侧为大气。求液体作用在侧壁下层部分上的净压力及其作用点深度，忽略界面张力。
\end{origexamplebox}
\begin{xesolution}
从两液体界面向下量 $x$，$0\le x\le h_2$。上层已贡献表压 $\rho_1gh_1$，下层再贡献 $\rho_2gx$，故
\[
 \Delta p(x)=\rho_1gh_1+\rho_2gx,
 \qquad \dd F=\Delta p(x)b\,\dd x.
\]
每一水平窄条所受净压力均沿侧壁法线向外，因而可以直接相加其大小：
\[
 F=bg\left(\rho_1h_1h_2+\frac12\rho_2h_2^2\right).
\]
求作用点不能直接取几何中心，因为下方压强更大。设合力位于界面下方 $x_c$，对界面处水平轴取矩，有
\[
 Fx_c=\int_0^{h_2}x\,\dd F
 =bg\left(\frac12\rho_1h_1h_2^2+\frac13\rho_2h_2^3\right),
\]
所以
\[
 x_c=\frac{\rho_1h_1h_2^2/2+\rho_2h_2^3/3}
 {\rho_1h_1h_2+\rho_2h_2^2/2}.
\]
距自由面的深度是 $h_1+x_c$。没有上层时 $x_c=2h_2/3$；若上层压力远大于下层压强变化，分布接近均匀，$x_c\to h_2/2$。两极限分别对应三角形和近似矩形压强分布，且 $\rho gbh^2$ 的量纲确为力。
\end{xesolution}

\subsection{浮力来自压力差，稳定性还取决于力矩}
设想把浸没物体占据的空间重新填满周围液体。静止时，这部分液体的重量被周围压力合力平衡；只要边界上的压强分布相同，把它换成物体后，压力合力也不变。这一替换论证给出浮力，而不是另外假设一种向上的力。
对于跨越两层液体的物体，应分别计算 $\rho_1gV_1+\rho_2gV_2$，不能把总排水体积都乘同一个密度。浮力与重力相等只保证竖直合力为零；物体稍微倾斜后是否产生回复力矩，还需比较重心与浮力作用线，不能仅据“浮着”判断稳定。

\begin{origexamplebox}{浮体的竖直小振动}\label{ch13:ex:teaching-buoyancy}
质量为 $m$、水平截面积恒为 $S$ 的直立柱体，在密度 $\rho$ 的液体中平衡漂浮。只允许竖直运动，忽略流体附加惯性、阻力和表面张力，且小振动中底面始终浸没、顶面始终露出。求平衡浸深及小振动周期。
\end{origexamplebox}
\begin{xesolution}
平衡时 $mg=\rho gSh_0$，所以 $h_0=m/(\rho S)$。取柱体向下的位移 $x$ 为正；此时浸深为 $h_0+x$，浮力向上且增大：
\[
 m\ddot x=mg-\rho gS(h_0+x)=-\rho gSx.
\]
平衡项消去后，剩下的正是指向平衡位置的回复力。因此
\[
 \omega^2=\frac{\rho gS}{m},\qquad
 T=2\pi\sqrt{\frac{m}{\rho gS}}.
\]
若 $x<0$，浮力减少，合力向下，也把柱体推回平衡。$\rho gS$ 的单位为 $\mathrm{N/m}$，相当于等效劲度系数。真实液体随物体一起加速会增加有效惯性，故本题结果依赖题设的忽略条件；也不能把它用于已经完全浸没、排液体积不随 $x$ 变化的柱体。
\end{xesolution}
